\startmode[single]

% \enabletrackers[publications, publications.crossref, publications.details, publications.cite, publications.strings]
%
% \usebtxdataset[main][bib/default.bib,bib/physik.bib,bib/mathe.bib,bib/seminarfach.bib]
% \usebtxdefinitions[aps]
% \setupbtx[dataset=main]
% \definebtxrendering[bibrendering][aps][dataset=main]

\setvariables[meta]
  [title={Kernphysik:},
   subtitle={Zerfallsarten / Strahlungsarten},
   date={2025-03-12},
   author={Niklas von Hirschfeld},
  ]

\setupinteraction[title={\getvariable{meta}{title}}]

\starttext

\startfrontmatter
\component c_titlepage
\component c_toc
\stopfrontmatter

\setuppagenumber[number=1]

\startbodymatter

\startchapter[title=\getvariable{meta}{subtitle}]

\stopmode

\section{Aufgabe 1}

In den dargestellten Versuchsbeobachtungen werden die Eigenschaften von $\alpha
-, \beta -, \gamma -$ Strahlungen getestet. Einmal in einem Magnetfeld (links)
und mit verschiedenen Materialien (rechts).

Zu beobachten ist, dass $\alpha -$ Strahlung in dem Magnetischen feld abgelenkt
wird und zwar in die entgegen gesetzte Richtung, wie die in die die $\beta -$
Strahlung abgelenkt wird. Sie muss also die entgegengesetze Ladung besitzen
($\alpha -$ Strahlung besitzt eine positive Ladun $+2e$ und $\beta -$Strahlung
eine negative $-e$).
$\gamma -$ Strahlung hingegen hat keine Ladung und wird nicht in dem Feld
abgelenkt. Auch kann man beobachten, dass $\alpha -$ Strahlung eine größere
Masse als $\beta -$ Strahlung, da sie nicht so stark abgelenkt wird.

Bei den Materialien wird beobachtet, dass $\alpha -$ nicht durch das Papier
durchdringen kann. $\beta -$ hingegen tritt durch das Papier durch, wird aber
von Alu geblock. $\gamma -$ Strahlung geht durh Papier und Alu hindurch und
wird von Blei abgeblockt.

\section{Aufgabe 2}

\startitemize[packed]
\item Var 1: Magnetfeld, Photoplatte unten, gucken, wie starkt abgelenkt
\item Var 2: Elektrisches Feld, ...
\item Var 3: Wien-Filter: senkrechtes Magnetfeld + elektrisches Querfled, so
einstellen, dass Teilchen exakt geradeaus fliegen.
\stopitemize

\startmode[single]

\stopchapter

\stopbodymatter

\startbackmatter
% \component c_definitions
\component c_bibliography
\stopbackmatter
\stoptext

\stopmode
